% Самостійний документ: xelatex (двічі).
% 60 хв викладу; усі запитання та пошук відповідей поза цим часом.
\documentclass[aspectratio=169,11pt]{beamer}
\usepackage{fontspec,polyglossia}
\setdefaultlanguage{ukrainian}
\setmainfont{DejaVu Serif}
\setsansfont{DejaVu Sans}
\setmonofont{DejaVu Sans Mono}
\newfontfamily\cyrillicfont{DejaVu Serif}
\newfontfamily\cyrillicfontsf{DejaVu Sans}
\newfontfamily\cyrillicfonttt{DejaVu Sans Mono}
\usepackage{amsmath,amssymb,booktabs,array}
\usefonttheme{professionalfonts}
\definecolor{ink}{HTML}{183341}
\definecolor{accent}{HTML}{087C83}
\setbeamercolor{normal text}{fg=ink,bg=white}
\setbeamercolor{structure}{fg=accent}
\setbeamercolor{frametitle}{fg=ink}
\setbeamertemplate{navigation symbols}{}
\setbeamertemplate{footline}{\hfill\insertframenumber\hspace{5mm}\vspace{3mm}}
\setbeamersize{text margin left=8mm,text margin right=8mm}
\setbeamerfont{frametitle}{size=\large,series=\bfseries}
\hypersetup{colorlinks=true,urlcolor=accent,linkcolor=ink}
\newcommand{\unit}[1]{\,\text{#1}}
\newcommand{\src}[2]{\par\vfill{\tiny\href{#1}{#2}}}
\newcommand{\E}{\operatorname{E}}
\newcommand{\Var}{\operatorname{Var}}
\newcommand{\xref}{x_{\mathrm{ref}}}
\title{Похибки і статистична\newline довіра до результату}
\subtitle{Лекція 4. Математичний конспект і практичні завдання}
\author{Метрологія, технологічні вимірювання та прилади}
\date{}
\begin{document}
\begin{frame}\titlepage\end{frame}

\begin{frame}{Як користуватися конспектом}
\begin{itemize}
\item Продовження лекції 3: канал сформував показ; тепер описуємо його відхилення та статистичну мінливість.
\item Чотири блоки: похибка і якість; розподіли; інтервали; статистичний звіт.
\item $X$ — випадкова величина; $x_i$ — отриманий показ; $\mu,\sigma$ — параметри моделі; $\bar x,s$ — вибіркові оцінки.
\item Числові серії навчальні, крім окремо названого прикладу NIST.
\item Наприкінці — чотири окремі слайди з розгорнутими Python-задачами без готових розв'язків.
\end{itemize}
60 хвилин відведено на виклад; запитання, пошук і обговорення відповідей не входять у цей час.
\end{frame}

\section{Похибка та показники якості}
\begin{frame}{Визначення: форми похибки}
\[
\Delta=x-\xref,\qquad
\delta=\frac{\Delta}{\xref}\cdot100\%,\quad \xref\ne0,
\qquad \gamma=\frac{\Delta}{x_n}\cdot100\%,\quad x_n>0.
\]
\begin{itemize}
\item $\Delta$ має одиницю величини та знак; «абсолютна» не означає $|\Delta|$.
\item $\xref$ — прийняте опорне значення; його власну невизначеність потрібно враховувати у повній метрологічній задачі.
\item $x_n$ — явно визначене нормувальне значення; у наших прикладах це ширина діапазону з нульовим початком.
\item При $\xref=0$ відносна похибка за цією формулою не визначена; у програмі повертаємо \texttt{None}.
\end{itemize}
Відсоткові приклади використовують додатне $\xref$; для температури у $^\circ$C відносну похибку тут не обчислюємо.
\src{https://jcgm.bipm.org/vim/en/2.16.html}{VIM 2.16: measurement error}
\end{frame}

\begin{frame}{Одна різниця — три способи подання}
Навчальні дані: $x=101\unit{кПа}$, $\xref=100\unit{кПа}$, $x_n=200\unit{кПа}$.
\[
\Delta=+1\unit{кПа},\qquad \delta=+1\%,\qquad \gamma=+0{,}5\%.
\]
\medskip
Для $x=99\unit{кПа}$ за тієї самої опори $\Delta=-1\unit{кПа}$: знак повідомляє напрям відхилення.
\medskip

Якщо два прилади мають умовну межу $1\%$ діапазону, то для діапазонів $[0;10]\unit{бар}$ і $[0;100]\unit{бар}$ абсолютні межі становлять відповідно $0{,}1\unit{бар}$ і $1\unit{бар}$.
\medskip

Менший відсоток при іншому знаменнику сам по собі не означає меншого фізичного відхилення.
\end{frame}

\begin{frame}{Визначення: складові та ознаки похибки}
Навчальна модель стабільного режиму відносно фіксованої опори:
\[
X=\xref+b+\varepsilon,\qquad \E[\varepsilon]=0,
\qquad b=\E[X]-\xref.
\]
\begin{itemize}
\item $b$ — зміщення відносно опори; оцінена поправка має протилежний знак, але не нульову невизначеність.
\item $\varepsilon$ — випадкове відхилення від середнього рівня; усереднення не усуває $b$.
\item Груба помилка запису або процедури потребує встановленої причини; незвичний показ ще не доводить помилки.
\item Статична/динамічна ознака описує режим: наприклад, запізнення при зміні входу може створювати динамічне відхилення.
\end{itemize}
Систематичний ефект у ширшій задачі може змінюватися передбачувано; стала $b$ — спрощення для цього режиму.
\end{frame}

\begin{frame}{Визначення: показники якості}
\small
\begin{tabular}{>{\raggedright\arraybackslash}p{0.23\textwidth}>{\raggedright\arraybackslash}p{0.69\textwidth}}
\toprule
Правильність & Близькість середнього нескінченної кількості повторень до опорного значення.\\[1mm]
Прецизійність & Близькість повторних показів між собою за визначених умов.\\[1mm]
Збіжність & Прецизійність за умов повторюваності.\\[1mm]
Відтворюваність & Прецизійність за визначених змінених умов, наприклад між лабораторіями.\\[1mm]
Точність & Якісне поняття близькості результату до істинного значення.\\[1mm]
Достовірність & Обґрунтованість довіри до висновку; довірчий рівень є лише одним статистичним аспектом.\\
\bottomrule
\end{tabular}
\src{https://jcgm.bipm.org/vim/en/2.13.html}{VIM 2.13–2.15, 2.21, 2.25; точність не є окремою числовою величиною}
\end{frame}

\section{Розподіли випадкових похибок}
\begin{frame}{Визначення: випадкова величина і щільність}
Для неперервної моделі зі щільністю $f$:
\[
f(x)\ge0,\qquad \int_{-\infty}^{\infty}f(x)\,dx=1,
\qquad P(a\le X\le b)=\int_a^b f(x)\,dx.
\]
\begin{itemize}
\item Щільність у точці не є ймовірністю цієї точки; ймовірність інтервалу дорівнює площі під кривою.
\item Для дискретного $X$: $P(X=x_j)=p_j$, $p_j\ge0$, $\sum_jp_j=1$.
\item Серія $x_1,\ldots,x_n$ — реалізації моделі, а не сама модель.
\item Тренд або змішування режимів можуть порушити припущення про однаковий розподіл усіх спостережень.
\end{itemize}
\end{frame}

\begin{frame}{Визначення: математичне сподівання і дисперсія}
За існування відповідних моментів:
\[
\mu=\E[X]=\int_{-\infty}^{\infty}x f(x)\,dx,\qquad
\sigma^2=\E[(X-\mu)^2],\qquad \sigma=\sqrt{\sigma^2}.
\]
У дискретній моделі інтеграл замінюється сумою: $\E[X]=\sum_jx_jp_j$.
\[
\E[X+c]=\mu+c,\qquad \Var(X+c)=\sigma^2,
\qquad \Var(aX)=a^2\sigma^2.
\]
Отже, спільне зміщення змінює центр, але не розсіювання.
\medskip

Якщо $X$ має одиницю кПа, то $\sigma$ також має одиницю кПа, а $\sigma^2$ — кПа$^2$.
\end{frame}

\begin{frame}{Визначення: нормальна модель}
\[
X\sim\mathcal N(\mu,\sigma^2),\quad \sigma>0,\qquad
f(x)=\frac{1}{\sigma\sqrt{2\pi}}
\exp\!\left[-\frac{(x-\mu)^2}{2\sigma^2}\right].
\]
\[
Z=\frac{X-\mu}{\sigma}\sim\mathcal N(0,1),\qquad
P(|Z|\le1)\approx0{,}6827,\quad P(|Z|\le2)\approx0{,}9545.
\]
\begin{itemize}
\item Хвости не обриваються: $\sigma$ не є найбільшим можливим відхиленням.
\item Модель багатьох малих внесків потребує умов; наявність шуму сама по собі не доводить нормальності.
\item Нормальність середніх великих незалежних вибірок за умов центральної граничної теореми не означає нормальності окремих показів.
\end{itemize}
\src{https://www.itl.nist.gov/div898/handbook/eda/section3/eda3661.htm}{NIST: Normal Distribution}
\end{frame}

\begin{frame}{Визначення: рівномірна модель}
\[
\varepsilon\sim U(-a;a),\quad a>0,\qquad
f(e)=\begin{cases}1/(2a),&-a\le e\le a,\\0,&\text{поза інтервалом}.\end{cases}
\]
\[
\E[\varepsilon]=0,\qquad
\Var(\varepsilon)=\frac1{2a}\int_{-a}^{a}e^2\,de
=\frac{a^2}{3},\qquad \sigma=\frac{a}{\sqrt3}.
\]
Для ідеального округлення до найближчого рівня з кроком $q$:
\[
a=q/2\quad\Longrightarrow\quad \sigma=\frac q{\sqrt{12}}.
\]
При $q=1\,{}^\circ\mathrm C$ отримуємо $\sigma\approx0{,}289\,{}^\circ\mathrm C$ для внеску округлення за рівномірної моделі.
\src{https://www.analog.com/media/en/training-seminars/tutorials/MT-001.pdf}{Analog Devices, MT-001: модель ідеального квантування}
\end{frame}

\begin{frame}{Коли модель потребує перегляду}
\begin{itemize}
\item Постійний вхід може давати постійний код: повторення не створюють незалежного рівномірного шуму квантування.
\item Сума двох незалежних $U(-a;a)$ має трикутну щільність на $[-2a;2a]$, а не рівномірну.
\item Перемикання режимів може створити дві групи; тренд означає зміну рівня з часом.
\item Коротка серія дозволяє обчислити описові показники, але не доводить вибраний закон розподілу.
\end{itemize}
Спочатку перевіряємо механізм, часовий порядок та умови, потім застосовуємо статистичну модель.
\end{frame}

\section{Точкові та інтервальні оцінки}
\begin{frame}{Визначення: точкові оцінки}
Для $n\ge2$ спостережень:
\[
\bar x=\frac1n\sum_{i=1}^n x_i,\qquad
s^2=\frac1{n-1}\sum_{i=1}^n(x_i-\bar x)^2,\qquad s=\sqrt{s^2}.
\]
\begin{itemize}
\item $\bar x$ оцінює $\mu$, $s^2$ оцінює $\sigma^2$, $s$ оцінює $\sigma$.
\item Для незалежних однаково розподілених спостережень зі скінченною дисперсією $\E[\bar X]=\mu$, $\E[S^2]=\sigma^2$.
\item Незміщеність $S^2$ не означає незміщеності $S=\sqrt{S^2}$.
\item У Python: \texttt{statistics.mean}, \texttt{variance}, \texttt{stdev}.
\end{itemize}
Позначення великими літерами стосується оцінювачів до отримання вибірки, малими — їхніх обчислених значень.
\end{frame}

\begin{frame}{Визначення: стандартне відхилення середнього}
Для незалежних однаково розподілених $X_i$:
\[
\Var(\bar X)=\frac1{n^2}\sum_{i=1}^n\Var(X_i)
=\frac{\sigma^2}{n},\qquad
\sigma_{\bar X}=\frac\sigma{\sqrt n},\qquad
\widehat{\sigma}_{\bar X}=\frac s{\sqrt n}.
\]
\begin{itemize}
\item $s$ — розсіювання показів; $s/\sqrt n$ (SEM) — оцінка розсіювання середніх серій.
\item За залежності виникають також коваріації:
\end{itemize}
\[
\Var(\bar X)=\frac1{n^2}\left(\sum_i\Var(X_i)+2\sum_{i<j}\operatorname{Cov}(X_i,X_j)\right).
\]
Коваріація описує спільну зміну відхилень; її ігнорування може занизити оцінку розсіювання середнього.
\end{frame}

\begin{frame}{Визначення: двобічний інтервал Стьюдента}
При $X_i\overset{\mathrm{iid}}\sim\mathcal N(\mu,\sigma^2)$, $\sigma>0$, $n\ge2$:
\[
T=\frac{\bar X-\mu}{S/\sqrt n}\sim t_{n-1},\qquad
P\!\left(-t_{1-\alpha/2,n-1}\le T\le t_{1-\alpha/2,n-1}\right)=1-\alpha.
\]
Звідси процедура інтервального оцінювання:
\[
h=t_{1-\alpha/2,n-1}\frac{s}{\sqrt n},\qquad
\mathrm{CI}_{1-\alpha}=[\bar x-h;\bar x+h].
\]
\textbf{iid} означає незалежні однаково розподілені спостереження. Для 95\%: $\alpha=0{,}05$, $1-\alpha/2=0{,}975$.
\src{https://www.itl.nist.gov/div898/handbook/eda/section3/eda352.htm}{NIST: Confidence Limits for the Mean}
\end{frame}

\begin{frame}{Критичний множник і ширина}
\begin{center}
\begin{tabular}{ccc}
\toprule
$n$ & $n-1$ & $t_{0{,}975,n-1}$\\\midrule
5 & 4 & 2,776\\
10 & 9 & 2,262\\
20 & 19 & 2,093\\
$\to\infty$ & $\to\infty$ & $\to1{,}960$\\\bottomrule
\end{tabular}
\end{center}
\[
h=t\,s/\sqrt n.
\]
За сталих $s$ і $t$ чотириразове збільшення $n$ вдвічі зменшує $h$; у малих серіях змінюється також $t$. Вищий довірчий рівень розширює інтервал.
\src{https://www.itl.nist.gov/div898/handbook/eda/section3/eda3672.htm}{NIST: Critical Values of the Student's t Distribution}
\end{frame}

\begin{frame}{Як тлумачити 95\%}
\begin{itemize}
\item При багаторазовому повторенні всієї процедури приблизно 95\% інтервалів накриють фіксований параметр $\mu$, якщо припущення виконані.
\item Для вже обчисленого інтервалу твердження $\mu\in\mathrm{CI}$ або істинне, або хибне; 95\% стосується процедури покриття.
\item Це не інтервал для 95\% окремих показів і не гарантія відсутності зміщення.
\end{itemize}
\textbf{Приклад NIST:} $n=195$, $\bar x=9{,}261460$, $s=0{,}022789$, $t\approx1{,}9723$; 95\% CI:
\[
[9{,}258242;9{,}264679].
\]
Фізичну одиницю цим числам тут не приписуємо. Інтервал для середнього показу не замінює повний бюджет невизначеності.
\src{https://www.itl.nist.gov/div898/handbook/eda/section3/eda352.htm}{NIST: вихідні числа та межі інтервалу}
\end{frame}

\section{Статистичні оцінки характеристик похибок}
\begin{frame}{Чому знаменник дорівнює n − 1}
\[
d_i=x_i-\bar x,\qquad \sum_i d_i=0.
\]
Один залишок визначається рештою. Алгебраїчно:
\[
\sum_i(X_i-\bar X)^2
=\sum_i(X_i-\mu)^2-n(\bar X-\mu)^2.
\]
Для незалежної однакової моделі зі скінченною дисперсією:
\[
\E\!\left[\sum_i(X_i-\bar X)^2\right]
=n\sigma^2-n\frac{\sigma^2}{n}=(n-1)\sigma^2.
\]
Тому ділення на $n-1$ дає незміщену оцінку дисперсії. Ділення на $n$ описує середній квадрат відхилень у самому заданому наборі.
\end{frame}

\begin{frame}{Визначення: оцінка зміщення й розсіювання}
Для фіксованої прийнятої опори:
\[
\Delta_i=x_i-\xref,\qquad
\widehat b=\bar\Delta=\bar x-\xref,
\qquad \Delta_i-\bar\Delta=x_i-\bar x.
\]
Отже, вибіркове СКВ похибок відносно їхнього середнього дорівнює СКВ показів.
\medskip

\textbf{Дві окремі відповіді:} $\widehat b$ оцінює зміщення відносно опори, $s$ описує випадкове розсіювання серії.
\medskip

Без опори можна порівнювати спостережувану повторюваність за однакових умов, але не доводити правильність. Невизначеність опори не зникає від збільшення числа повторів робочого приладу.
\end{frame}

\begin{frame}{Аномалія — привід перевірити умови}
\begin{itemize}
\item Навчальна серія: $10{,}0;10{,}1;9{,}9;10{,}0;15{,}0\unit{кПа}$.
\item Перевіряємо первинний запис, одиниці, час, перемикання режимів та стан обладнання.
\item Доведена помилка може бути виправлена або виключена з документуванням причини.
\item Невідома причина не дає права тихо видалити показ; реальна подія може вимагати іншої моделі.
\item Формальні тести на викиди також мають припущення; автоматичну відсіч у цій лекції не застосовуємо.
\end{itemize}
Початковий набір зберігаємо; рішення про склад оброблюваної серії робимо перевірним.
\src{https://www.itl.nist.gov/div898/handbook/eda/section3/eda35h.htm}{NIST: Detection of Outliers}
\end{frame}

\begin{frame}{Подання результату}
Звіт: величина, одиниця, умови, $n$, $\bar x$, $s$, опора та $\widehat b$ за наявності, тип і рівень інтервалу.
\medskip

Навчальне правило компактного подання: дві значущі цифри в $h$, середнє — до того самого десяткового розряду.
\[
\bar x=12{,}345678\unit{кПа},\quad h=0{,}047891\unit{кПа}
\]
\[
\Longrightarrow\quad\bar x=12{,}346\unit{кПа},\quad h=0{,}048\unit{кПа}.
\]
Підпис: «95\% CI для середнього $\mu$: $\bar x\pm h$; калібрувальні та інші внески тут не включено».
\medskip

Обчислюємо за початковими значеннями; округлюємо тільки подання. Конкретна лабораторна методика може встановлювати інші правила оформлення.
\end{frame}

\begin{frame}{Покриття офіційних питань}
\small
\begin{tabular}{cp{0.80\textwidth}}
\toprule
20 & Різновиди похибок вимірювань.\\[1mm]
21 & Форми відображення кількісних характеристик похибок вимірювань.\\[1mm]
22 & Показники якості вимірювань.\\[1mm]
23 & Закони розподілу похибок і результатів вимірювань.\\[1mm]
24 & Точкові і інтервальні ймовірності оцінки випадкових похибок вимірювань.\\[1mm]
25 & Статистичні оцінки характеристик похибок вимірювань.\\
\bottomrule
\end{tabular}
\medskip

Блок 1: 20–22; блок 2: 23; блок 3: 24; блок 4: 25. Формулювання звірено з офіційним DOCX 2020 року; у питанні 24 збережено текст оригіналу.
\end{frame}

\begin{frame}{Джерела та засоби для самостійної роботи}
\small
\begin{itemize}
\item \href{https://jcgm.bipm.org/vim/en/2.16.html}{JCGM, VIM: похибка}; \href{https://jcgm.bipm.org/vim/en/2.13.html}{точність}; \href{https://jcgm.bipm.org/vim/en/2.15.html}{прецизійність}.
\item \href{https://www.itl.nist.gov/div898/handbook/eda/section3/eda3661.htm}{NIST: нормальна модель}; \href{https://www.itl.nist.gov/div898/handbook/eda/section3/eda3662.htm}{рівномірна модель}.
\item \href{https://www.analog.com/media/en/training-seminars/tutorials/MT-001.pdf}{Analog Devices, Walt Kester, MT-001: ідеальне квантування}.
\item \href{https://www.itl.nist.gov/div898/handbook/eda/section3/eda352.htm}{NIST: інтервал для середнього}; \href{https://www.itl.nist.gov/div898/handbook/eda/section3/eda3672.htm}{критичні значення Стьюдента}; \href{https://www.itl.nist.gov/div898/handbook/eda/section3/eda35h.htm}{аномалії}.
\item \href{https://docs.python.org/3/library/statistics.html}{Python: statistics}; \href{https://docs.python.org/3/library/math.html}{Python: math}.
\end{itemize}
Наступні чотири слайди — окремі розгорнуті завдання. Потрібні лише базовий Python і стандартні бібліотеки; готові розв'язки не наведено.
\end{frame}

\appendix
\begin{frame}{Практичне завдання 1: порівняння з опорою}
\small
\textbf{Вхід:} \texttt{[49.8, 50.1, 50.2, 49.9]} кПа; $\xref=50$ кПа, $x_n=100$ кПа.
\begin{enumerate}
\item Напишіть функцію, що приймає список, опору і нормувальне значення; для кожного показу повертає $x_i,\Delta_i,\delta_i,\gamma_i$.
\item Додайте середнє $\Delta_i$ через \texttt{statistics.mean}; поясніть, що це оцінка зміщення відносно прийнятої опори.
\item При $\xref=0$ поверніть \texttt{None} для $\delta_i$; при $x_n\le0$ або порожньому списку — зрозуміле повідомлення.
\item Виведіть таблицю з одиницями; збережіть знаки; початковий список не змінюйте.
\end{enumerate}
\textbf{Перевірки:} показ дорівнює опорі; покази по різні боки від опори; нульова опора; недодатне $x_n$; порожній список.
\medskip

\textbf{Подати:} функцію, результати перевірок і два речення про відмінність абсолютної, відносної та приведеної форм; повну якість приладу за ними не оголошувати.
\end{frame}

\begin{frame}{Практичне завдання 2: паспорт двох серій}
\small
\textbf{Вхід, $^\circ$C:}\par
\texttt{A = [20.0, 20.1, 19.9, 20.0, 20.0]}\par
\texttt{B = [19.6, 20.4, 19.8, 20.2, 20.0]}.
\begin{enumerate}
\item Створіть функцію, що повертає $n,\bar x,s^2,s$ через \texttt{mean}, \texttt{variance}, \texttt{stdev}.
\item Підпишіть одиниці: $^\circ$C для середнього й СКВ, $({}^\circ\mathrm C)^2$ для дисперсії; поясніть знаменник $n-1$.
\item Порівняйте центр і розсіювання окремо; не робіть висновку про правильність без опори чи про нормальність без перевірки моделі.
\item Для $n<2$ повідомте про недостатність даних; не підміняйте вибіркові функції \texttt{pvariance}/\texttt{pstdev}.
\end{enumerate}
\textbf{Перевірки:} однакові числа; один елемент; порожній список; додавання сталої до кожного показу.
\medskip

\textbf{Подати:} функцію, два звіти й коротке пояснення того, що зберігається при спільному зсуві всіх значень.
\end{frame}

\begin{frame}{Практичне завдання 3: інтервал для середнього}
\small
\textbf{Вхід:} \texttt{[100.1, 99.9, 100.2, 100.0, 99.8]} кПа.\par
Модель: незалежні нормальні спостереження, невідома $\sigma$; 95\%, \texttt{tcrit = 2.776} для $n=5$.
\begin{enumerate}
\item Функція приймає список і \texttt{tcrit}, повертає $n,\bar x,s,s/\sqrt n$, нижню та верхню межі.
\item Використайте \texttt{statistics.mean}, \texttt{statistics.stdev}, \texttt{math.sqrt}; не пишіть обчислення квантиля.
\item Вимагайте $n\ge2$ та додатного \texttt{tcrit}; округлюйте лише звіт; зазначте, що інтервал стосується $\mu$ за заданої моделі.
\item При зміні довжини списку узгодьте множник із $n-1$ та рівнем довіри; таблиця наведена раніше.
\end{enumerate}
\textbf{Перевірки:} симетричність меж відносно середнього; збільшення множника; додавання сталої до всіх показів; список із одного числа.
\medskip

\textbf{Подати:} функцію, звіт і одне речення про те, чому цей CI не охоплює автоматично систематичні та калібрувальні внески.
\end{frame}

\begin{frame}{Практичне завдання 4: перевірний лабораторний звіт}
\small
\textbf{Вхід, В:}\par
\texttt{A = [5.01, 5.00, 4.99, 5.02, 4.98]}\par
\texttt{B = [5.01, 5.00, 4.99, 5.02, 5.60]}; $\xref=5{,}00$ В.
\begin{enumerate}
\item Для A прийміть незалежну нормальну модель; виведіть $n,\bar x,s^2,s,\widehat b$ та 95\% CI із \texttt{tcrit = 2.776}; повторно використайте функцію із задачі 3.
\item Для B причина 5.60 невідома: збережіть усі дані, подайте описові оцінки та повідомлення про перевірку журналу, без автоматичної відсічі й заявленого 95\% CI.
\item Відхиляйте порожні/надто короткі списки, текст, \texttt{None} та нескінченні значення; для чисел можна використати \texttt{math.isfinite}.
\item Форматуйте величини у В до трьох знаків, дисперсію у В$^2$ — до шести; проміжні значення не округлюйте.
\end{enumerate}
\textbf{Подати:} функції, два звіти, перевірку незмінності списків і два речення про розсіювання та необхідні додаткові перевірки.
\end{frame}
\end{document}
